% GENERATED FILE. Edit canonical lessons, then run npm run generate:books.
% canonical-source-hash: fnv1a64:311cc8b0958ad983
% canonical-lessons: RU-C292-obyskivat, RU-C292-arestovyvat, RU-C292-doprashivat, RU-C292-dopros, RU-C292-pokazaniya

\chapter{Search and questioning}
\label{ch:ru-search-and-questioning}
\hlchaptermodality{292}

\begin{chapteropening}
\textbf{By the end of this chapter:} \emph{I can say that the police searched and arrested someone, and talk about questioning and statements.}

Complete the last lesson of chapter 292: I can say that the police searched and arrested someone, and talk about questioning and statements.
\end{chapteropening}

\section[\texorpdfstring{\ru{обыскивать} (obýskivat)}{обыскивать (obýskivat)}]{\ru{обыскивать} (obýskivat) — to search (a person or a place)}

\label{lesson:RU-C292-obyskivat}



\begin{quote}
Before the new one: say the Russian for a suspect, then the Russian for to suspect.
\end{quote}

\subsection*{You'll want to know: \ru{обыскивать}}
\textbf{\ru{обыскивать}} (\emph{obýskivat}) — \textquotedblleft{}to search (a person or a place)\textquotedblright{}.

\begin{quote}
\ru{Полицейские} \ru{обыскали} \ru{квартиру}.
\end{quote}

\begin{quote}
\emph{Politséyskiye obyskáli kvartíru.}
\end{quote}

\begin{quote}
The police searched the flat.
\end{quote}

Imperfective; perfective \emph{\ru{обыскать}} (\emph{obyskát}), \emph{\ru{обыщу}}. \emph{\ru{Искать}} is to look for; \emph{\ru{обыскивать}} goes through a place. \emph{\ru{Обыск}} — a search.

\begin{cousinweb}
\emph{\ru{обыскивать}} is built on \emph{\ru{об}-} \textquotedblleft{}all round\textquotedblright{} and \emph{\ru{искать}} \textquotedblleft{}to look for\textquotedblright{}, which is inherited from Proto-Slavic \emph{jĭskati}, from PIE \emph{h\textsubscript{2}eys-} \textquotedblleft{}to wish, to seek\textquotedblright{} (EtymDB).
\end{cousinweb}

\begin{grammarlens}[title={What you've built}]
One more word for telling, explaining and describing at B1.
\end{grammarlens}

\subsection*{Your turn}
\begin{itemize}
\raggedright
  \item \emph{Say it:} \emph{obýskivat}
  \item \emph{Say it:} \emph{obýskivat}, once more
  \item \emph{Say it:} a sentence of your own with \emph{obýskivat}
  \item \emph{Recall:} say the Russian for a suspect, then the Russian for to suspect, then say \emph{obýskivat} again
\end{itemize}

\begin{tcolorbox}[breakable,colback=teal!4,colframe=teal!35!black,title={Before you move on}]
What does \ru{обыскивать} mean? (\textquotedblleft{}To search (a person or a place)\textquotedblright{}.) Say it once more.
\end{tcolorbox}
\section[\texorpdfstring{\ru{арестовывать} (arestóvyvat)}{арестовывать (arestóvyvat)}]{\ru{арестовывать} (arestóvyvat) — to arrest}

\label{lesson:RU-C292-arestovyvat}



\begin{quote}
Before the new one: say the Russian for to suspect, then the Russian for to search.
\end{quote}

\subsection*{You'll want to know: \ru{арестовывать}}
\textbf{\ru{арестовывать}} (\emph{arestóvyvat}) — \textquotedblleft{}to arrest\textquotedblright{}.

\begin{quote}
\ru{Его} \ru{арестовали} \ru{в} \ru{аэропорту}.
\end{quote}

\begin{quote}
\emph{Yevó arestováli v aeroportú.}
\end{quote}

\begin{quote}
He was arrested at the airport.
\end{quote}

Imperfective; perfective \emph{\ru{арестовать}} (\emph{arestovát}). \emph{\ru{Арест}} — an arrest.

\begin{cousinweb}
\emph{\ru{арестовывать}} is built on \emph{\ru{арест}} \textquotedblleft{}an arrest\textquotedblright{}, which is usually traced through German \emph{Arrest} to Medieval Latin \emph{arrestum} \textquotedblleft{}a stopping\textquotedblright{}; EtymDB lists no parent.
\end{cousinweb}

\begin{grammarlens}[title={What you've built}]
One more word for telling, explaining and describing at B1.
\end{grammarlens}

\subsection*{Your turn}
\begin{itemize}
\raggedright
  \item \emph{Say it:} \emph{arestóvyvat}
  \item \emph{Say it:} \emph{arestóvyvat}, once more
  \item \emph{Say it:} a sentence of your own with \emph{arestóvyvat}
  \item \emph{Recall:} say the Russian for to suspect, then the Russian for to search, then say \emph{arestóvyvat} again
\end{itemize}

\begin{tcolorbox}[breakable,colback=teal!4,colframe=teal!35!black,title={Before you move on}]
What does \ru{арестовывать} mean? (\textquotedblleft{}To arrest\textquotedblright{}.) Say it once more.
\end{tcolorbox}
\section[\texorpdfstring{\ru{допрашивать} (dopráshivat)}{допрашивать (dopráshivat)}]{\ru{допрашивать} (dopráshivat) — to question, to interrogate}

\label{lesson:RU-C292-doprashivat}



\begin{quote}
Before the new one: say the Russian for to search, then the Russian for to arrest.
\end{quote}

\subsection*{You'll want to know: \ru{допрашивать}}
\textbf{\ru{допрашивать}} (\emph{dopráshivat}) — \textquotedblleft{}to question, to interrogate\textquotedblright{}.

\begin{quote}
\ru{Следователь} \ru{допрашивал} \ru{свидетелей} \ru{весь} \ru{день}.
\end{quote}

\begin{quote}
\emph{Slédovatel dopráshival svidételey ves den.}
\end{quote}

\begin{quote}
The investigator questioned witnesses all day.
\end{quote}

Imperfective; perfective \emph{\ru{допросить}} (\emph{doprosít}), \emph{\ru{допрошу}}. It takes the accusative.

\begin{cousinweb}
\emph{\ru{допрашивать}} is built on \emph{\ru{до}-} and the root of \emph{\ru{просить}} \textquotedblleft{}to ask\textquotedblright{}, inherited from Proto-Slavic \emph{prositi}, from PIE \emph{pre\'{k}-} \textquotedblleft{}to ask\textquotedblright{} (EtymDB): to ask to the end.
\end{cousinweb}

\begin{grammarlens}[title={What you've built}]
One more word for telling, explaining and describing at B1.
\end{grammarlens}

\subsection*{Your turn}
\begin{itemize}
\raggedright
  \item \emph{Say it:} \emph{dopráshivat}
  \item \emph{Say it:} \emph{dopráshivat}, once more
  \item \emph{Say it:} a sentence of your own with \emph{dopráshivat}
  \item \emph{Recall:} say the Russian for to search, then the Russian for to arrest, then say \emph{dopráshivat} again
\end{itemize}

\begin{tcolorbox}[breakable,colback=teal!4,colframe=teal!35!black,title={Before you move on}]
What does \ru{допрашивать} mean? (\textquotedblleft{}To question, to interrogate\textquotedblright{}.) Say it once more.
\end{tcolorbox}
\section[\texorpdfstring{\ru{допрос} (doprós)}{допрос (doprós)}]{\ru{допрос} (doprós) — questioning, an interrogation}

\label{lesson:RU-C292-dopros}



\begin{quote}
Before the new one: say the Russian for to arrest, then the Russian for to question.
\end{quote}

\subsection*{You'll want to know: \ru{допрос}}
\textbf{\ru{допрос}} (\emph{doprós}) — \textquotedblleft{}questioning, an interrogation\textquotedblright{}.

\begin{quote}
\ru{Допрос} \ru{длился} \ru{три} \ru{часа}.
\end{quote}

\begin{quote}
\emph{Doprós dlílsya tri chasá.}
\end{quote}

\begin{quote}
The questioning lasted three hours.
\end{quote}

Masculine. \emph{\ru{На} \ru{допросе}} — under questioning.

\begin{cousinweb}
\emph{\ru{допрос}} is built on \emph{\ru{до}-} and the root of \emph{\ru{просить}} \textquotedblleft{}to ask\textquotedblright{}, Proto-Slavic \emph{prositi}, from PIE \emph{pre\'{k}-} (EtymDB).
\end{cousinweb}

\begin{grammarlens}[title={What you've built}]
One more word for telling, explaining and describing at B1.
\end{grammarlens}

\subsection*{Your turn}
\begin{itemize}
\raggedright
  \item \emph{Say it:} \emph{doprós}
  \item \emph{Say it:} \emph{doprós}, once more
  \item \emph{Say it:} a sentence of your own with \emph{doprós}
  \item \emph{Recall:} say the Russian for to arrest, then the Russian for to question, then say \emph{doprós} again
\end{itemize}

\begin{tcolorbox}[breakable,colback=teal!4,colframe=teal!35!black,title={Before you move on}]
What does \ru{допрос} mean? (\textquotedblleft{}Questioning, an interrogation\textquotedblright{}.) Say it once more.
\end{tcolorbox}
\section[\texorpdfstring{\ru{показания} (pokazániya)}{показания (pokazániya)}]{\ru{показания} (pokazániya) — a statement, testimony}

\label{lesson:RU-C292-pokazaniya}



\begin{quote}
Before the new one: say the Russian for to question, then the Russian for questioning.
\end{quote}

\subsection*{You'll want to know: \ru{показания}}
\textbf{\ru{показания}} (\emph{pokazániya}) — \textquotedblleft{}a statement, testimony\textquotedblright{}.

\begin{quote}
\ru{Свидетель} \ru{дал} \ru{показания}.
\end{quote}

\begin{quote}
\emph{Svidétel dal pokazániya.}
\end{quote}

\begin{quote}
The witness gave a statement.
\end{quote}

Usually plural in this sense. \emph{\ru{Давать} \ru{показания}} — to testify. The same plural also means the readings of a meter.

\begin{cousinweb}
\emph{\ru{показания}} is built on \emph{\ru{показать}} \textquotedblleft{}to show\textquotedblright{}, from \emph{\ru{казать}}, inherited from Proto-Slavic \emph{kazati} (EtymDB): what a witness shows.
\end{cousinweb}

\begin{grammarlens}[title={What you've built}]
One more word for telling, explaining and describing at B1.
\end{grammarlens}

\subsection*{Your turn}
\begin{itemize}
\raggedright
  \item \emph{Say it:} \emph{pokazániya}
  \item \emph{Say it:} \emph{pokazániya}, once more
  \item \emph{Say it:} a sentence of your own with \emph{pokazániya}
  \item \emph{Recall:} say the Russian for to question, then the Russian for questioning, then say \emph{pokazániya} again
\end{itemize}

\begin{tcolorbox}[breakable,colback=teal!4,colframe=teal!35!black,title={Before you move on}]
What does \ru{показания} mean? (\textquotedblleft{}A statement, testimony\textquotedblright{}.) Say it once more.
\end{tcolorbox}
